% Chapter 10: The Practitioner's Paradigm: Infinite-Dimensional Data on GPUs
% Part II of the book.
% STATUS: book-standard prose (Part II).

\chapter{The Practitioner's Paradigm: Infinite-Dimensional Data on GPUs}
\label{ch:part2-the-practitioner-s-paradigm-infinite-dimen}
\label{sec:practitioner}
% -----------------------------------------------------------------------------

Part~I is written about functions; a training loop runs on arrays. Nothing in the
preceding nine chapters closes that distance, and the gap is where implementations fail
for reasons unrelated to approximation theory.
\index{graphics processing unit (GPU)|textbf}A \emph{graphics processing unit} (GPU)
holds no more of a field than any other memory does: $B \times C \times N_1 \times
\cdots \times N_d$ floating-point numbers, finite along every axis. The
``infinite-dimensional'' of this chapter's title names the object those numbers
approximate, not the object they are. What separates an operator-learning program from
an image-classification program is a discipline rather than a data type. Shape, loss,
and rescaling are each chosen so that the procedure tracks a map between function spaces
instead of a map between arrays of one privileged size.
\chapref{ch:part1-introduction-the-paradigm-shift-to-infinit} ended by promising this
translation explicitly---from $\mathcal{A}$, $\mathcal{U}$, and $N$ grid points to the
axes $(B, C, H, W)$---and the first section below supplies it.

The anchor is Definition~\ref{def:ch06-discretization-invariance} of
\chapref{ch:part1-discretization-invariance-and-continuum-li}, which states what an
implementation owes the continuum: a parameter space that does not grow with the node
count, a map defined on function spaces, and a discrete realization whose output
converges to that map as the mesh is refined, in the sense of \eqref{eq:limit_commute}.
Four engineering decisions answer it, one per section below: how a sampled field is laid
out in memory, how the network is told where it
is, how error is measured, and how fields are rescaled before the network sees them.
Each is the discrete shadow of a Part~I object---the restriction operator $P_N$, the
arguments of the kernel in \eqref{eq:fredholm_kernel}, the population risk
\eqref{eq:ch01-risk}, and the input measure $\mu \in \mathcal{P}(\mathcal{A})$---and each
has a failure mode: a number that looks reasonable on the training grid and means nothing
off it.

Prerequisites are practical: a \texttt{torch.nn.Module} written in
PyTorch~\cite{paszke2019pytorch}, the ability to read a batched tensor contraction, and
the vocabulary of device computing---\textsc{cuda} (Compute Unified Device Architecture),
the platform on which kernels are launched onto the GPU, and \textsc{cuDNN} (CUDA Deep
Neural Network library), the convolution and transform routines built on top of it. One
numerical-analysis habit matters more than the rest: whenever a program sums over grid
nodes, ask what integral the sum approximates and with what weights, since that decides
whether a printed number is a property of the field or a property of the mesh. No result
here is new; what follows is the dictionary, with the cost and the failure behavior of
each entry.

% -----------------------------------------------------------------------------
\section{Tensor Layouts for Function Spaces}
\label{sec:ch21-layouts}

In computer vision, an image is a discrete grid $X \in \mathbb{R}^{B \times C \times H
\times W}$, and the array is the whole of the object. In operator learning the same
tensor stands for a continuous function $a: \Omega \to \mathbb{R}^{C}$ sampled on a
spatial discretization $\mathcal{M}_N = \{x_j\}_{j=1}^{N}$: it is the output of the
restriction operator $P_N$ of Definition~\ref{def:ch06-discretization-invariance}, one
row per node, stacked for $B$ independent draws from the input measure. Two layouts
dominate production workflows, and both should be held in mind for the rest of
Part~II. \index{tensor layout|textbf}\index{data layout!grid versus point cloud}

The first is the structured grid layout, or channel-first layout, used when
$\Omega = [0, 1]^d$ is a uniform Cartesian mesh and the sample is the array
\begin{equation}
  \label{eq:ch21-grid-layout}
  X \in \mathbb{R}^{B \times C \times N_1 \times \dots \times N_d},
\end{equation}
where $B$ is the batch size, $C$ is the number of field components---the velocity pair
$(u, v)$ together with the pressure $p$ gives $C = 3$---and the $N_j$ are the spatial
grid dimensions. The translation \chapref{ch:part1-introduction-the-paradigm-shift-to-infinit}
promised is then exact, apart from the batch axis: $C$ is the channel count $d_a$ of
Part~I, the product $\prod_j N_j$ is the node count $N$, and the slice
$X[b, \cdot, n_1, \dots, n_d]$ is the sample $a(x_j)$ at the node whose multi-index is
$(n_1, \dots, n_d)$. Running the realized operator produces
the same shape with $C = d_u$. That excluded axis is the only one carrying no operator
semantics at all: the members of a batch are exchangeable draws, so $B$ is chosen at
inference time by memory and nothing else, which is why \eqref{eq:ch21-grid-layout} can
describe a thousand permeability fields and one field equally well.
\index{batch axis}\index{channel axis}

Storage cost is arithmetic, and it is the constraint that shapes every other decision.
A field held in single precision (float32, four bytes per real number) occupies
\begin{equation}
  \label{eq:ch21-memory}
  \mathrm{bytes}_{\mathrm{real}} = 4\, B\, C \prod_{j=1}^{d} N_j,
  \qquad
  \mathrm{bytes}_{\mathrm{spec}} = 8\, B\, C\, N_1 \big(\lfloor N_2/2 \rfloor + 1\big),
\end{equation}
the second expression being the complex-valued spectrum a two-dimensional real transform
returns (complex64, two float32 per entry, conjugate symmetry discarding half the
frequencies). Numerically, a batch of $B = 8$ at $C = 32$ costs about $4.2$~MB on a
$64 \times 64$ grid and $67$~MB on a $256 \times 256$ one, with the transform
intermediate a further $4.3$~MB and $68$~MB respectively. Refining fourfold along each
axis multiplies the node count, and with it the activation memory, by sixteen, while the
parameter tensor does not move at all.
That asymmetry explains why the architectures of this book hold the mode and channel
budgets fixed across a training run and pay for a finer grid out of the batch axis; it is
read off \eqref{eq:ch21-memory} and is the reason the mode count of
\chapref{ch:part1-the-fourier-neural-operator-spectral-theor} is a design parameter
rather than a consequence of the data.
\index{memory footprint!of a grid tensor}\index{float32|textbf}

What the grid layout buys is a kernel library, and the purchase is worth itemizing.
Because the spatial axes are trailing and contiguous, a batched real \emph{fast Fourier
transform} (FFT) over the last two axes (\texttt{torch.fft.rfft2}, the two-axis
real-to-complex transform) is one call rather
than a gather, and channel mixing at each retained frequency is a single Einstein-summation
contraction (\texttt{torch.einsum}) of the form used in
\chapref{ch:part2-fourier-neural-operators-fno-from-scratch}~\cite{li2021fourier};
\textsc{cuDNN} two- and
three-dimensional convolutions consume the same tensor without transposition.
\index{fast Fourier transform (FFT)!batched, on-device}\index{rfft2 (real two-dimensional FFT)}\index{einsum (Einstein summation)}
The second layout is the point-cloud or mesh layout, used for irregular geometries such
as aerodynamic surfaces and arterial trees, where a function is represented as a set of
nodes and their coordinates,
\begin{equation}
  \label{eq:ch21-node-layout}
  X_{\text{nodes}} \in \mathbb{R}^{B \times N_{\text{nodes}} \times C}, \qquad
  X_{\text{coords}} \in \mathbb{R}^{B \times N_{\text{nodes}} \times d}.
\end{equation}
\index{point-cloud layout}\index{unstructured mesh!node-array layout}
The channel axis has moved to the end because the elementary operation there is pointwise
in the node index: the lift $\mathcal{P}$ in the composition \eqref{eq:no_composition} is
a shallow multilayer perceptron (MLP) acting on the channel dimensions, so placing $C$
innermost makes one matrix multiplication serve every node. The node axis is
also, in general, of different length for different samples, which a rectangular grid
tensor cannot express.

The two layouts differ along a single axis, and naming it settles the comparison: a
weight tensor in \eqref{eq:ch21-grid-layout} is indexed by lattice offsets and frequency
bins, whereas a weight tensor in \eqref{eq:ch21-node-layout} is indexed by the physical
displacement $x_i - x_j$. That is precisely the distinction drawn in
\eqref{eq:ch06-cnn}, where the discrete stencil of a convolutional neural network (CNN)
shrinks in physical space as the mesh refines, and in \eqref{eq:ch06-kernel-discrete},
where a kernel is rebuilt from a continuous function on each new mesh. The grid layout
inherits translation equivariance and the $\mathcal{O}(N \log N)$ cost of
\eqref{eq:ch05-convolution-theorem}; the mesh layout inherits neither, aggregating over
neighbor lists at a cost set by the number of edges, and it must carry the quadrature
weights $w_j$---Voronoi dual areas---explicitly, since there is no constant spacing to
absorb them. Drop those weights and the sum approximates nothing, which is the
failure \chapref{ch:part1-discretization-invariance-and-continuum-li} warns against in
the strongest terms; the graph implementations of
\chapref{ch:part2-graph-neural-operators-gno-and-irregular-g} exist to keep
them attached~\cite{li2023geometry}.
\index{quadrature weights!carried explicitly on meshes}

Two further cautions close the comparison. Converting between the layouts is not free:
grid to nodes is a plain gather when the nodes coincide with grid points and an
interpolation when they do not, while nodes to grid is an interpolation whose error was
never seen during training---the objection
\chapref{ch:part1-introduction-the-paradigm-shift-to-infinit} raises against
re-interpolation, and which
\chapref{ch:part1-discretization-invariance-and-continuum-li} prices under the name
``interpolate, then run''. And a layout error is not
always detected by a shape check, because a transform will oblige whichever axes it is
handed, returning a tensor of exactly the expected size whose contents are meaningless.
The check is semantic, not dimensional: the axes being transformed must be the axes along
which the operator is nonlocal.

% -----------------------------------------------------------------------------
\section{Positional Grid Encoding}
\label{sec:ch21-grid-embedding}

A critical empirical discovery in neural operator engineering is that passing only the
field values $a(x)$ often leads to translational ambiguity. To break the symmetry and
supply geometric context, practitioners explicitly append normalized spatial coordinates
as additional input channels,
\begin{equation}
  \label{eq:ch21-augment}
  a_{\text{augmented}}(x) = \big[ a(x), \, x_1, \, x_2, \, \dots, \, x_d \big] \in \mathbb{R}^{C_{\text{in}} + d}.
\end{equation}
\index{positional grid encoding|textbf}\index{coordinate channels!as input augmentation}
The concatenation runs along the channel axis, so the pointwise lift $\mathcal{P} :
\mathbb{R}^{d_a} \to \mathbb{R}^{d_v}$ of \eqref{eq:no_composition} must instead be sized
for $C_{\text{in}} + d$ inputs; the extra $d$ channels
are not data but the argument of the function, broadcast to the shape of the data.

The mechanism behind the empirical report is architectural rather than statistical, and
for the periodic case it is exact. A shift of the input by $s$ multiplies each Fourier
coefficient by a phase $e^{-2\pi i k \cdot s}$, and the multiplier $R_\phi(k)$ of
\eqref{eq:ch05-multiplier} acts before that phase can be compared with anything, so the
spectral branch commutes with translation; the local bypass $W_\ell$ and the pointwise
activation in \eqref{eq:fno_layer} do the same. A network built from those pieces
cannot express a preference for absolute location, so any target that is not
translation-invariant---one with a boundary layer, a localized forcing, a Dirichlet
condition---must be reached through something that knows where it is. The density results
of \chapref{ch:part1-operator-approximation-theory-and-universa} do not require the
coordinates, and nothing forbids them either; what they are is the cheapest available way
of returning a little inhomogeneity to an architecture whose premise, as
\chapref{ch:part1-the-fourier-neural-operator-spectral-theor} puts it, is a homogeneous
medium. \index{translation ambiguity!broken by coordinate channels}

The price is a loss of regularity, and Part~I already has its number. A coordinate ramp
is not periodic: it jumps across the boundary of $[0,1]^d$ when the grid is understood as
a torus, which is exactly what the discrete transform assumes.
\chapref{ch:part1-discretization-invariance-and-continuum-li} records the consequence
for the discretization term---augmented inputs carry regularity $H^{1/2 - \epsilon}$ for
every $\epsilon > 0$ and no better, and because the exponent in \eqref{eq:ch06-rate} is
the regularity of the least smooth object in the pipeline, the observed decay under
refinement never reaches $N^{-1}$. The ramp is thus bought against convergence under
refinement, not against accuracy at fixed
resolution; the two are different ledgers and conflating them is a common source of
puzzlement in ablation studies. \index{coordinate channels!regularity cost}

\begin{lstlisting}[language=Python, caption={Appending the two normalized coordinate planes as extra channels of a grid field: input $(B, C, H, W)$, output $(B, C{+}2, H, W)$.}, label={lst:ch21-grid-embedding}]
import torch

class GridEmbeddingND(torch.nn.Module):
    """Appends continuous spatial coordinates as extra channels to input fields."""
    def __init__(self, dim: int = 2):
        super().__init__()
        self.dim = dim

    def forward(self, x: torch.Tensor) -> torch.Tensor:
        # Input shape: (B, C, H, W)
        batch_size, _, height, width = x.shape
        grid_x = torch.linspace(0, 1, steps=height, device=x.device, dtype=x.dtype)
        grid_y = torch.linspace(0, 1, steps=width, device=x.device, dtype=x.dtype)
        grid_x, grid_y = torch.meshgrid(grid_x, grid_y, indexing='ij')
        
        # Expand across batch dimension: (B, 2, H, W)
        grid = torch.stack([grid_x, grid_y], dim=0).unsqueeze(0).repeat(batch_size, 1, 1, 1)
        return torch.cat([x, grid], dim=1) # Shape: (B, C + 2, H, W)
\end{lstlisting}

Listing~\ref{lst:ch21-grid-embedding} is the utility that performs \eqref{eq:ch21-augment}
in production, and three things should be taken from it. First, the shape choreography:
input $(B, C, H, W)$, two one-dimensional ramps of lengths $H$ and $W$, a meshgrid with
\texttt{indexing='ij'} giving two $(H, W)$ planes, stacking to $(2, H, W)$, a leading
singleton and a repeat giving $(B, 2, H, W)$, and concatenation on dimension 1 giving the
$(B, C + 2, H, W)$ the caption of the listing names. Second, the grid is constructed on
\texttt{x.device} with \texttt{x.dtype}, so the coordinates never force a host-to-device
copy and never silently up-cast a half-precision tensor---a point that
\chapref{ch:part2-training-regimes-optimization-and-failure-} takes up under automatic
mixed precision (AMP). Third, the module is named for general $d$ but its \texttt{forward}
implements the two-dimensional case, and its stored \texttt{dim} is not consulted there;
a three-dimensional user must not trust the name. Note also that the ramps run from $0$
to $1$ inclusive, so the last node receives $1$ rather than $1 - \Delta x$: the jump
discussed above is present in the data by construction, and the memory cost of
\texttt{repeat} is $2BHW$ elements even though every batch element receives the same
planes. \index{grid embedding!shape choreography}

On the point-cloud side of \eqref{eq:ch21-node-layout} no such augmentation is needed:
the coordinates already exist as a separate tensor, and the kernels of
\chapref{ch:part2-graph-neural-operators-gno-and-irregular-g} consume them as arguments
rather than as channels. The axis of difference is worth stating once: grid encoding
injects geometry into the feature vector, where the lift can mix it with the field,
whereas mesh layouts keep geometry in the index structure, where it selects which nodes
interact. Which is preferable is a question about the discretization, not about the
architecture.

% -----------------------------------------------------------------------------
\section{Discretization-Invariant Loss Functions}
\label{sec:ch21-loss}

Standard mean squared error (\texttt{torch.nn.MSELoss}) sums squared differences across
array elements, $\text{MSE} = \frac{1}{N} \sum_i (u_i - \hat{u}_i)^2$. Under mesh
refinement ($N \to 4N$) the magnitude of that quantity fluctuates with the array
convention: reduced by summation it grows with the node count, and reduced by a mean it
settles toward a Riemann sum of the squared error but still carries the field's own
scale, so a batch containing one large-magnitude solution dominates the objective.
Operator learning therefore uses the \textbf{relative $L^2$ norm} as an approximation of
the continuous functional metric,
\begin{equation}
  \label{eq:rel_l2_loss}
  \mathcal{L}_{\text{rel-}L^2}(u, \hat{u}) = \frac{\| u - \hat{u} \|_{L^2(\Omega)}}{\| u \|_{L^2(\Omega)}} \approx \frac{\left( \sum_{i=1}^N |u_i - \hat{u}_i|^2 w_i \right)^{1/2}}{\left( \sum_{i=1}^N |u_i|^2 w_i \right)^{1/2}}.
\end{equation}
For equispaced grids, the quadrature weights $w_i = \Delta x$ cancel out between numerator
and denominator, rendering \eqref{eq:rel_l2_loss} strictly invariant to the grid density
$N$. \index{relative $L^2$ error (metric)|textbf}

Read \eqref{eq:rel_l2_loss} against the population risk \eqref{eq:ch01-risk} of
\chapref{ch:part1-introduction-the-paradigm-shift-to-infinit} and its place in the
argument is clear: the display is the inner quotient of that expectation, evaluated as a
quadrature on one sample, and the training loss adds the outer average over the batch.
Two conventions differ and both appear in the literature. The risk squares the ratio;
the implementation below takes the root and then averages, which preserves the ranking of
samples by error but not their weighting during optimization, since
the derivative of a square root is largest where the error is smallest. The percentages
reported for the benchmark problems of
\chapref{ch:part2-canonical-pde-benchmarks-and-reference-wor} and measured in
Appendix~\ref{app:b} are of this root-then-average form. What the invariance does
\emph{not} supply is a mesh-independent error value: the cancellation makes the
\emph{metric} scale-free, while the number computed still changes with $N$ because $u$
and $\hat{u}$ are different functions sampled on different grids, by the amount the
discretization term \eqref{eq:ch06-rate} describes. \index{loss function!discretization-invariant}

\begin{lstlisting}[language=Python, caption={Per-sample relative $\ell^p$ error reduced over every axis but the batch: with \texttt{p=2} this evaluates the quotient of the relative $L^2$ metric, and it accepts differentiated tensors for the $H^1$ variant.}, label={lst:ch21-lploss}]
class LpLoss(torch.nn.Module):
    """Computes relative L^p loss across arbitrary spatial resolutions."""
    def __init__(self, p: int = 2, reduction: str = 'mean', eps: float = 1e-7):
        super().__init__()
        self.p = p
        self.reduction = reduction
        self.eps = eps

    def forward(self, pred: torch.Tensor, target: torch.Tensor) -> torch.Tensor:
        # Shapes: (B, C, ...)
        diff = pred - target
        # Flatten spatial and channel dimensions
        diff_norm = torch.norm(diff.reshape(diff.shape[0], -1), p=self.p, dim=1)
        target_norm = torch.norm(target.reshape(target.shape[0], -1), p=self.p, dim=1)
        
        rel_error = diff_norm / (target_norm + self.eps)
        if self.reduction == 'mean':
            return torch.mean(rel_error)
        elif self.reduction == 'sum':
            return torch.sum(rel_error)
        return rel_error
\end{lstlisting}

Listing~\ref{lst:ch21-lploss} is what a training loop actually calls, and the reader
should learn from it exactly which parts of \eqref{eq:rel_l2_loss} the code keeps. Both
norms are taken after flattening every axis but the batch, so the discrete $\ell^p$ norm
is computed over the concatenation of channels and nodes---the vector-valued
$L^2(\Omega; \mathbb{R}^{C})$ norm of the notation table, with one number per sample---and
the ratio is formed per sample before reduction. Weights $w_i$ do not appear anywhere,
which is correct precisely because constant weights cancel in the quotient, and is wrong
the moment the nodes are not equispaced: applied to a point-cloud tensor from
\eqref{eq:ch21-node-layout}, the same call returns an unweighted sum that is not the
$L^2$ metric of anything. The guard \texttt{eps} sits in the denominator only, which
prevents division by zero but does nothing for conditioning when $\|u\|$ is genuinely
small---a nearly constant or nearly zero label---a case that recurs in the
derivative-normalized terms of \eqref{eq:ch07-sobolev-loss} and that a robust workflow
handles by normalizing over a batch instead of per sample.
\index{loss function!relative $L^2$ versus mean squared error}
\index{quadrature weights!cancellation in relative losses}

The caption of the listing names the $H^1$ Sobolev loss as well, and the relation is
worth spelling out because the code contains no derivatives. With \texttt{p=2} the call
returns \eqref{eq:rel_l2_loss}; each derivative-normalized term of
\eqref{eq:ch07-sobolev-loss} is that same quotient taken on differentiated fields, and a
single inhomogeneous $H^1$ quotient
$\|u - \hat{u}\|_{H^1}/\|u\|_{H^1}$ is obtained by stacking $[u, \partial_1 u, \dots,
\partial_d u]$ as extra channels of the same call, or by weighting modal errors by
$(1 + |k|^2)^{1/2}$ inside the transform. Either route
costs one extra derivative of prediction \emph{and} label---spectral differentiation on a
grid, a difference operator on a mesh---and buys control over the quantity most
applications actually read, the flux, the vorticity, the wall shear stress. \index{$H^1$ loss!derivative terms supplied by the caller}
Finally, the reduction argument: \texttt{'none'} is exposed because a single averaged
scalar hides the sample-wise distribution that distinguishes a surrogate with a few
catastrophic failures from one with uniformly mediocre error, and because the
per-sample values are what the resolution studies of
\chapref{ch:part2-canonical-pde-benchmarks-and-reference-wor} report. The precision of
these reductions, and of the long summations a relative denominator entails, is taken up
with the rest of the training budget in
\chapref{ch:part2-training-regimes-optimization-and-failure-}.

% -----------------------------------------------------------------------------
\section{Infinite-Dimensional Data Normalization}
\label{sec:ch21-normalization}

Standard batch normalization or scalar min-max scaling disrupts the spatial correlation
structure of continuous fields, because each averages over the wrong set. Batch
normalization keeps a parameter per channel but estimates its statistics by pooling the
batch and, in its spatial form, the positions as well, so what it produces is not the
spread of the field at a location but the spread over locations; the numbers it reports
also depend on which samples happened to share a step, and since its parameters carry no
spatial axis, refining the grid changes what they pool over without giving them a location
to depend on. Scalar min-max folds every position
and every channel into one pair of numbers, discarding the fact that variability in a
solution field is strongly position-dependent---large near an inlet or an interface,
small where the boundary condition pins the value. Practitioners therefore use
\textbf{pointwise unit-Gaussian normalizers} computed across the training ensemble,
\begin{equation}
  \label{eq:ch21-normalize}
  \mu(x) = \frac{1}{M} \sum_{j=1}^M u^{(j)}(x), \qquad \sigma(x) = \sqrt{\frac{1}{M} \sum_{j=1}^M \big(u^{(j)}(x) - \mu(x)\big)^2 + \epsilon}.
\end{equation}
\index{data normalization!pointwise (per-location)}
\index{batch normalization!versus pointwise normalization}

The axis along which the three schemes differ is thus the set being averaged over:
positions, the whole array, or the ensemble. Only the last produces objects that are
functions of $x$. In the grid layout of \eqref{eq:ch21-grid-layout}, $\mu$ and $\sigma$
carry a channel index as well and have shape $(C, N_1, \dots, N_d)$---one field pair, not
$B$ of them---so they are computed once offline and reused for the lifetime of the
model. Their relation to Part~I is direct: the ensemble $\{u^{(j)}\}_{j=1}^{M}$ is a
sample from the input measure $\mu \in \mathcal{P}(\mathcal{A})$ under which
\eqref{eq:ch01-risk} takes its expectation, and \eqref{eq:ch21-normalize} is a Monte
Carlo estimate of that measure's pointwise mean and standard deviation. The two letters
collide and should be disambiguated at once: $\mu(x)$ in \eqref{eq:ch21-normalize} is a
function on $\Omega$, whereas the $\mu$ of \eqref{eq:ch01-risk} is a measure on function
space, the object the former estimates. Two consequences
follow from reading the display that way. The divisor $\epsilon$ sits inside the square
root, so where the sample variance is negligible the normalizer saturates at
$\sqrt{\epsilon}$ rather than amplifying round-off---the locations where an
ensemble-estimated standard deviation is least reliable are exactly those where it is
smallest. And the estimator inherits the support and the size of the ensemble: like the
surrogate itself it licenses no claim outside the training distribution, and the
denominator $1/M$ rather than $1/(M-1)$ is the population convention, a difference
immaterial at large $M$ and visible when the ensemble is deliberately small.
\index{input measure!empirical mean and standard deviation}
\index{standardization!ensemble statistics}

During inference on a refined grid $N' > N$, the mean $\mu$ and standard deviation
$\sigma$ fields are interpolated to the new coordinates, preserving continuous
normalization. This is the sentence that most deserves its qualification, because it is
where the resolution-transfer property is weakest and where it is easiest to overclaim.
The repair works if $\mu$ and $\sigma$ are smooth enough for the interpolation to be
accurate, and that smoothness is an assumption about the measure, not a gift of
averaging: an ensemble mean of rough fields is still rough. Where the input measure
concentrates on low-regularity data, interpolating the normalizer injects an
interpolation error that the training loop never observed, in the same way that
re-interpolating a surrogate's output does---the objection
\chapref{ch:part1-introduction-the-paradigm-shift-to-infinit} raises against
re-interpolation. The honest statement of the situation is that
\eqref{eq:ch21-normalize} defines a pair of continuum fields, and that the code stores
discretizations of them, so the mesh-free claim belongs to the object and not to the
array. \index{normalizer!interpolation under refinement}

One interaction with Section~\ref{sec:ch21-loss} decides what a reported percentage
means. Normalization is not invisible to a relative metric: dividing by a
position-dependent $\sigma(x)$ is not a scalar rescaling, so the relative $L^2$ error of
the normalized pair and the relative $L^2$ error of the physical pair are different
numbers, and neither is bounded by the other without a bound on $\sigma$ above and below.
Workflows therefore differ by whether the loss is evaluated before or after the
denormalization step, and the benchmark results quoted in
\chapref{ch:part2-canonical-pde-benchmarks-and-reference-wor} and measured in
Appendix~\ref{app:b} are in physical units, which is the convention this book adopts when
accuracy is to be compared against a solver.
\index{denormalization!effect on reported error}

\takeaway{From Arrays Back to Operators}{Each decision of this chapter is the discrete
realization of one Part~I object: the layout of \eqref{eq:ch21-grid-layout} and
\eqref{eq:ch21-node-layout} is the restriction $P_N$; the encoding
\eqref{eq:ch21-augment} restores the location dependence the stationary kernel discards;
the metric \eqref{eq:rel_l2_loss} is the quotient inside the risk \eqref{eq:ch01-risk};
and the normalizer \eqref{eq:ch21-normalize} estimates the input measure. None of them
makes a tensor infinite-dimensional. Each keeps the array honest about the function it
stands for.}

\medskip
The four decisions above share one premise: input and output are discretized on the same
mesh, so a sampled field and its solution can be compared node by node.
\chapref{ch:part1-operator-approximation-theory-and-universa} treats that premise as an
artifact. Its integral layer \eqref{eq:fredholm_kernel} reads the input as a function and
returns values at arbitrary points, and the architecture built on the corresponding
idea---DeepONet, a Deep Operator Network---splits the two sides of
\eqref{eq:ch21-grid-layout} into two different tensors and two different
networks~\cite{lu2021learning}: a branch that consumes evaluations at $m$ fixed sensors
and a trunk that consumes continuous query coordinates $y$. Everything assembled here
must then be re-decided. Which layout carries the sensor vector, whether positional
encoding helps a network that already receives coordinates as its second input, how the
relative metric of \eqref{eq:rel_l2_loss} is formed when the label is an unordered set of
point values rather than a grid function, and which normalization applies to a tensor
with no spatial axes at all: these are the questions of
\chapref{ch:part2-deeponet-engineering-branch-trunk-networks}, where the branch-trunk
model is built in PyTorch and the sensor choices that Part~I's sampling-error bounds
make consequential are put into code.
